\section{Introduction}

Image manipulation detection is increasingly evaluated as a supervised prediction task: a model assigns an authenticity score and, in localization settings, produces a pixel-level map. Such measurements are necessary for comparing systems, but their interpretation depends on the evaluation data and operating conditions. Dataset-specific acquisition and editing traces, annotation conventions, and post-processing can all affect the measured result. For forensic use, a high score on one benchmark should therefore not be read as evidence of stable performance on another dataset or after an image has been transformed.

This concern is visible in the broader image-forensics literature. Methods have progressed from analyses of acquisition and compression traces to learned image-level detectors and localization systems, including approaches based on self-consistency, rich manipulation features, and diffusion priors \cite{farid2009imageforgery,verdoliva2020mediaforensics,huh2018fighting,zhou2018rich,wu2019mantra,yu2024diffforensics}. Newer benchmarks broaden the range of synthetic-image and manipulation sources \cite{ojha2023universal,yan2024df40,zhu2023genimage,du2025forensichub}. Yet benchmark scores are not directly interchangeable when dataset composition, reference-mask meaning, model thresholds, or image processing differ. A reliable evaluation should report those conditions alongside the scores and should retain enough provenance to verify which inputs and model outputs produced each reported result.

Two further questions remain relevant after an evaluation has been run. First, how sharply do detection and localization results change across independent datasets and controlled post-processing? Second, can the reported model outputs be regenerated from the same source files and frozen model artifacts in a separate process? Standard detection metrics answer neither question by themselves. Forensic guidance emphasizes documented methods, observations and limitations, while benchmark initiatives emphasize performance under varied operational conditions \cite{swgde2025imageauth,nist2025guardians,nistopenmfc}. In this study, we treat evidence auditability narrowly as record completeness, input-integrity verification and deterministic output replay. These properties support technical inspection; they do not establish causal attribution or legal admissibility.

We present ForenX-Bench as an auditable evaluation protocol and report its completed instantiation using frozen TruFor and MVSS-Net models on Columbia Uncompressed, corrected CASIA v1 and corrected COVERAGE. The study evaluates image-level detection and mask-based localization under a fixed threshold, applies JPEG recompression and resize--restore stress transformations, and records model inputs and derived outputs in machine-readable manifests. A seeded, manipulation-stratified sample is then rerun in an independent process with input hashes checked and numerical agreement assessed at a predeclared tolerance. We ask whether performance is consistent across these datasets and transformations, where mask-based comparisons require annotation-specific qualifications, and whether the retained outputs can be reproduced under the recorded execution conditions.

The contribution is an empirical, reproducible audit of existing detectors rather than a new detection architecture. Specifically, we (i) report cross-dataset detection and localization results for two frozen models across three manipulation datasets; (ii) quantify changes under controlled JPEG and resize--restore transformations without retuning the decision threshold; and (iii) provide an auditable evidence record with a stratified independent-process replay assessment. The scope is limited to the evaluated datasets, checkpoints, transformations and execution environment. A dual-branch detector pilot is retained as a supplementary negative result and is not presented as evidence of an algorithmic improvement.
